\def\ProblemCode{SEGLAZY}
\def\ProblemTitle{Segment Tree (Lazy propagation)}
\makeproblem{\ProblemCode}{\ProblemTitle}

Cho dãy số nguyên $a_1, a_2, \dots, a_n$.
Xử lý $q$ truy vấn thuộc một trong các dạng sau:

\begin{itemize}
    \item \boxed{1\ l\ r\ k}: Tăng mỗi phần tử trong đoạn $[l, r]$ lên $k$ đơn vị $(|k| \le 10^{9})$.
    \item \boxed{2\ l\ r}: Tính tổng các phần tử trong đoạn $[l, r]$.
    \item \boxed{3\ l\ r}: Tìm phần tử có giá trị nhỏ nhất trong đoạn $[l, r]$.
    \item \boxed{4\ l\ r}: Tìm phần tử có giá trị lớn nhất trong đoạn $[l, r]$.
\end{itemize}

% ===== INPUT DESCRIPTION =====
\inputDesc{}

\begin{itemize}
    \item Dòng đầu tiên gồm hai số nguyên dương $n, q\ (1 \le n, q \le 10^{5})$.
    \item Dòng thứ hai gồm $n$ số nguyên $a_1, a_2, \dots, a_n\ (|a_i| \le 10^{9})$.
    \item $q$ dòng tiếp theo, mỗi dòng chứa một truy vấn.
\end{itemize}

% ===== OUTPUT DESCRIPTION =====
\outputDesc{}

\begin{itemize}
    \item Với mỗi truy vấn loại 2, 3 hoặc 4, in ra một dòng chứa kết quả tương ứng.
\end{itemize}

% ===== SUBTASKS =====
% \noindent{\sffamily \textbf{Subtasks}}
% \begin{itemize}
%     \item \textbf{Subtask ?} ($?\%$ số điểm): Không có ràng buộc gì thêm.
% \end{itemize}

\vspace{0.5em}

\ExampleBox{1}{
5 9\\
-3 5 -2 4 -1\\
2 1 5\\
3 1 5\\
4 1 5\\
1 2 4 3\\
2 1 5\\
3 2 4\\
4 2 4\\
1 1 5 -2\\
2 1 3
}{
3\\
-3\\
5\\
12\\
1\\
8\\
0
}{}
